\section*{第20章：递归类型}
\addcontentsline{toc}{section}{第20章：递归类型}

\begin{taplsolutionentrybox}
\phantomsection
\label{sol:translator:20.2.2}
\textbf{练习20.2.2}

要证明的是\cref{fig:20-1}中的同构递归演算满足以下两项性质。

\textbf{进展性。}若闭项 \(t\) 良类型，则 \(t\) 要么是值，要么存在
\(t'\) 使 \(t\trans t'\)。证明对类型推导作归纳。旧句法形式沿用简单类型
lambda 演算的证明，只需处理两种新增项。

若最后一条规则是 \textsc{T-Fold}，则项形如 \(\mathsf{fold}[U]t_1\)。由
归纳假设，\(t_1\) 要么可以求值一步，此时用 \textsc{E-Fold} 推进；要么已经
是值 \(v_1\)，此时 \(\mathsf{fold}[U]v_1\) 本身就是值。

若最后一条规则是 \textsc{T-Unfold}，则项形如
\(\mathsf{unfold}[U]t_1\)。由归纳假设，若 \(t_1\) 可以求值一步，就用
\textsc{E-Unfold}。若 \(t_1\) 已经是值，还需使用新增的典范形式事实：
闭的 \(U=\mu X.T\) 型值只能形如 \(\mathsf{fold}[U]v\)。于是
\textsc{E-UnfldFld} 可用，整个项可以继续求值。

\textbf{保持性。}若 \(\Gamma\vdash t:T\) 且 \(t\trans t'\)，则
\(\Gamma\vdash t':T\)。证明对单步求值推导作归纳。两个同余规则直接使用
归纳假设，再分别应用 \textsc{T-Fold} 或 \textsc{T-Unfold}。

关键情形是
\[
\mathsf{unfold}[S](\mathsf{fold}[U]v_1)\trans v_1.
\]
因为原项良类型，\textsc{T-Unfold} 给出 \(S=\mu X.T_1\)，并且
\(\mathsf{fold}[U]v_1:S\)。再反演 \textsc{T-Fold}，可知折叠标注必须是
同一个递归类型 \(U=S\)，且
\[
\Gamma\vdash v_1:\taplsubst{X}{S}{T_1}.
\]
这正是原项经 \textsc{T-Unfold} 得到的结果类型，因此规约后的 \(v_1\)
保持原类型。

证明中最值得注意的两点，是为递归类型补充典范形式引理，以及在
\textsc{E-UnfldFld} 情形中借助良类型性证明两个运行时未比较的类型标注实际
相同。

\taplsolutionbacklink{ex:20.2.2}{返回练习20.2.2}
\end{taplsolutionentrybox}
